\section{Resultados e discussão dos resultados}\label{sec:resultados}

Esta seção apresenta os resultados obtidos nos experimentos de verificação biométrica com os modelos DeepPrint e FLARE (e suas respectivas variações de alinhamento e realce) sobre os conjuntos de dados das edições FVC2000, FVC2002 e FVC2004, considerando as partições de avaliação principal (Parte A) dos sensores reais (DB1 a DB3).

A Tabela~\ref{tab:eer_base} apresenta a comparação de EER obtida pelas arquiteturas base (sem realce) no ponto de mínima diferença entre FAR e FRR. As Tabelas~\ref{tab:eer_deepprint_enh} e~\ref{tab:eer_flare_enh} detalham o efeito das etapas de realce (PriorEnh e UNetEnh) sobre o desempenho do DeepPrint e do FLARE (alinhamento por regressão e votação), respectivamente.

% Tabelas unificadas: Modelos Base, DeepPrint com Realce e FLARE com Realce (Parte A, DB1-DB3)

\begin{table}[htbp]
\centering
\small
\caption{EER (\%) das arquiteturas base (sem realce) nos subconjuntos Parte A do FVC.}\label{tab:eer_base}
\setlength{\tabcolsep}{8pt}
\begin{tabular}{ll rrr}
\toprule
\textbf{Edição} & \textbf{Base} & \textbf{DeepPrint} & \textbf{FLARE (Reg.)} & \textbf{FLARE (Vot.)} \\
\midrule
FVC2000 & DB1-A & 3.36 & 0.81 & 0.81 \\
 & DB2-A & 2.69 & 0.54 & 0.19 \\
 & DB3-A & 3.85 & 0.63 & 0.59 \\
\cmidrule{2-5}
\multicolumn{2}{l}{\textbf{Média FVC2000}} & \textbf{3.30} & \textbf{0.66} & \textbf{0.53} \\
\midrule
FVC2002 & DB1-A & 5.20 & 0.44 & 0.25 \\
 & DB2-A & 8.19 & 0.62 & 0.69 \\
 & DB3-A & 8.51 & 3.32 & 2.11 \\
\cmidrule{2-5}
\multicolumn{2}{l}{\textbf{Média FVC2002}} & \textbf{7.30} & \textbf{1.46} & \textbf{1.02} \\
\midrule
FVC2004 & DB1-A & 17.05 & 0.25 & 0.24 \\
 & DB2-A & 5.27 & 1.44 & 0.88 \\
 & DB3-A & 4.91 & 0.38 & 0.31 \\
\cmidrule{2-5}
\multicolumn{2}{l}{\textbf{Média FVC2004}} & \textbf{9.07} & \textbf{0.69} & \textbf{0.48} \\
\midrule[\heavyrulewidth]
\multicolumn{2}{l}{\textbf{Média Geral}} & \textbf{6.56} & \textbf{0.94} & \textbf{0.67} \\
\bottomrule
\end{tabular}
\end{table}

\begin{table}[htbp]
\centering
\small
\caption{EER (\%) do DeepPrint com métodos de realce nos subconjuntos Parte A do FVC.}\label{tab:eer_deepprint_enh}
\setlength{\tabcolsep}{8pt}
\begin{tabular}{ll rrr}
\toprule
\textbf{Edição} & \textbf{Base} & \textbf{Original} & \textbf{PriorEnh} & \textbf{UNetEnh} \\
\midrule
FVC2000 & DB1-A & 3.36 & 12.45 & 4.66 \\
 & DB2-A & 2.69 & 3.41 & 2.88 \\
 & DB3-A & 3.85 & 3.51 & 3.78 \\
\cmidrule{2-5}
\multicolumn{2}{l}{\textbf{Média FVC2000}} & \textbf{3.30} & \textbf{6.46} & \textbf{3.77} \\
\midrule
FVC2002 & DB1-A & 5.20 & 5.10 & 5.29 \\
 & DB2-A & 8.19 & 7.15 & 6.84 \\
 & DB3-A & 8.51 & 9.70 & 8.66 \\
\cmidrule{2-5}
\multicolumn{2}{l}{\textbf{Média FVC2002}} & \textbf{7.30} & \textbf{7.32} & \textbf{6.93} \\
\midrule
FVC2004 & DB1-A & 17.05 & 16.69 & 17.00 \\
 & DB2-A & 5.27 & 5.10 & 5.27 \\
 & DB3-A & 4.91 & 5.01 & 5.54 \\
\cmidrule{2-5}
\multicolumn{2}{l}{\textbf{Média FVC2004}} & \textbf{9.07} & \textbf{8.93} & \textbf{9.27} \\
\midrule[\heavyrulewidth]
\multicolumn{2}{l}{\textbf{Média Geral}} & \textbf{6.56} & \textbf{7.57} & \textbf{6.66} \\
\bottomrule
\end{tabular}
\end{table}

\begin{table}[htbp]
\centering
\small
\caption{EER (\%) do FLARE com métodos de realce nos subconjuntos Parte A do FVC.}\label{tab:eer_flare_enh}
\setlength{\tabcolsep}{5pt}
\begin{tabular}{ll rrr rrr}
\toprule
 & & \multicolumn{3}{c}{\textbf{FLARE (Regressão)}} & \multicolumn{3}{c}{\textbf{FLARE (Votação)}} \\
\cmidrule(lr){3-5} \cmidrule(lr){6-8}
\textbf{Edição} & \textbf{Base} & \textbf{Base} & \textbf{Prior} & \textbf{UNet} & \textbf{Base} & \textbf{Prior} & \textbf{UNet} \\
\midrule
FVC2000 & DB1-A & 0.81 & 16.81 & 1.06 & 0.81 & 16.83 & 0.87 \\
 & DB2-A & 0.54 & 17.42 & 0.37 & 0.19 & 17.43 & 0.19 \\
 & DB3-A & 0.63 & 22.36 & 0.87 & 0.59 & 22.36 & 0.64 \\
\cmidrule{2-8}
\multicolumn{2}{l}{\textbf{Média FVC2000}} & \textbf{0.66} & \textbf{18.86} & \textbf{0.77} & \textbf{0.53} & \textbf{18.87} & \textbf{0.57} \\
\midrule
FVC2002 & DB1-A & 0.44 & 25.54 & 0.68 & 0.25 & 25.54 & 0.25 \\
 & DB2-A & 0.62 & 34.25 & 0.56 & 0.69 & 34.24 & 0.31 \\
 & DB3-A & 3.32 & 27.65 & 3.00 & 2.11 & 27.65 & 2.12 \\
\cmidrule{2-8}
\multicolumn{2}{l}{\textbf{Média FVC2002}} & \textbf{1.46} & \textbf{29.15} & \textbf{1.42} & \textbf{1.02} & \textbf{29.14} & \textbf{0.89} \\
\midrule
FVC2004 & DB1-A & 0.25 & 29.46 & 0.25 & 0.24 & 29.51 & 0.31 \\
 & DB2-A & 1.44 & 20.29 & 1.48 & 0.88 & 20.29 & 1.31 \\
 & DB3-A & 0.38 & 18.01 & 1.26 & 0.31 & 18.02 & 1.26 \\
\cmidrule{2-8}
\multicolumn{2}{l}{\textbf{Média FVC2004}} & \textbf{0.69} & \textbf{22.59} & \textbf{1.00} & \textbf{0.48} & \textbf{22.61} & \textbf{0.96} \\
\midrule[\heavyrulewidth]
\multicolumn{2}{l}{\textbf{Média Geral}} & \textbf{0.94} & \textbf{23.53} & \textbf{1.06} & \textbf{0.67} & \textbf{23.54} & \textbf{0.81} \\
\bottomrule
\end{tabular}
\end{table}


A Figura~\ref{fig:histograms_fvc2000} apresenta a comparação da distribuição dos escores de similaridade entre o modelo DeepPrint e o modelo FLARE (com alinhamento por votação) para os conjuntos DB1-A a DB4-A da edição FVC2000. Observa-se que no FLARE a distribuição dos escores de impostores (vermelho) permanece concentrada em valores baixos (próximos a 0,2), resultando em uma separação quase total em relação aos usuários genuínos (verde) e justificando as taxas de EER substancialmente menores em comparação ao DeepPrint.

\begin{figure}[htbp]
\centering
\subfloat[DeepPrint - DB1-A]{\includegraphics[width=0.48\textwidth]{images/charts/histograms/deep_print_fvc_2000_db1_a.png}}
\hfill
\subfloat[FLARE (Votação) - DB1-A]{\includegraphics[width=0.48\textwidth]{images/charts/histograms/flare_voting_pose_fvc_2000_db1_a.png}}

\subfloat[DeepPrint - DB2-A]{\includegraphics[width=0.48\textwidth]{images/charts/histograms/deep_print_fvc_2000_db2_a.png}}
\hfill
\subfloat[FLARE (Votação) - DB2-A]{\includegraphics[width=0.48\textwidth]{images/charts/histograms/flare_voting_pose_fvc_2000_db2_a.png}}

\subfloat[DeepPrint - DB3-A]{\includegraphics[width=0.48\textwidth]{images/charts/histograms/deep_print_fvc_2000_db3_a.png}}
\hfill
\subfloat[FLARE (Votação) - DB3-A]{\includegraphics[width=0.48\textwidth]{images/charts/histograms/flare_voting_pose_fvc_2000_db3_a.png}}

\subfloat[DeepPrint - DB4-A]{\includegraphics[width=0.48\textwidth]{images/charts/histograms/deep_print_fvc_2000_db4_a.png}}
\hfill
\subfloat[FLARE (Votação) - DB4-A]{\includegraphics[width=0.48\textwidth]{images/charts/histograms/flare_voting_pose_fvc_2000_db4_a.png}}

\caption{Comparação da distribuição dos escores de similaridade entre DeepPrint (coluna à esquerda) e FLARE com alinhamento por votação (coluna à direita) para os subconjuntos DB1-A a DB4-A do FVC2000.}\label{fig:histograms_fvc2000}
\end{figure}

Para avaliar o impacto dos métodos de realce sobre a capacidade discriminativa das representações, a Figura~\ref{fig:histograms_enhancement} apresenta as distribuições de escores obtidas com o emprego das redes PriorEnh (coluna à esquerda) e UNetEnh (coluna à direita) para os modelos DeepPrint e FLARE (com alinhamento por votação) nos subconjuntos DB1-A e DB2-A do FVC2000. Observa-se que, enquanto a UNetEnh preserva a separação nítida entre genuínos e impostores em ambos os modelos, a aplicação da PriorEnh provoca um deslocamento substancial da distribuição de genuínos em direção a escores menores, gerando forte sobreposição com os impostores e explicando a degradação observada nas taxas de erro.

\begin{figure}[htbp]
\centering
\subfloat[DeepPrint + PriorEnh - DB1-A]{\includegraphics[width=0.48\textwidth]{images/charts/histograms/deep_print_priorenh_fvc_2000_db1_a.png}}
\hfill
\subfloat[DeepPrint + UNetEnh - DB1-A]{\includegraphics[width=0.48\textwidth]{images/charts/histograms/deep_print_unetenh_fvc_2000_db1_a.png}}

\subfloat[DeepPrint + PriorEnh - DB2-A]{\includegraphics[width=0.48\textwidth]{images/charts/histograms/deep_print_priorenh_fvc_2000_db2_a.png}}
\hfill
\subfloat[DeepPrint + UNetEnh - DB2-A]{\includegraphics[width=0.48\textwidth]{images/charts/histograms/deep_print_unetenh_fvc_2000_db2_a.png}}

\subfloat[FLARE (Votação) + PriorEnh - DB1-A]{\includegraphics[width=0.48\textwidth]{images/charts/histograms/flare_priorenh_voting_pose_fvc_2000_db1_a.png}}
\hfill
\subfloat[FLARE (Votação) + UNetEnh - DB1-A]{\includegraphics[width=0.48\textwidth]{images/charts/histograms/flare_unetenh_voting_pose_fvc_2000_db1_a.png}}

\subfloat[FLARE (Votação) + PriorEnh - DB2-A]{\includegraphics[width=0.48\textwidth]{images/charts/histograms/flare_priorenh_voting_pose_fvc_2000_db2_a.png}}
\hfill
\subfloat[FLARE (Votação) + UNetEnh - DB2-A]{\includegraphics[width=0.48\textwidth]{images/charts/histograms/flare_unetenh_voting_pose_fvc_2000_db2_a.png}}

\caption{Comparação da distribuição dos escores obtidos com os métodos de realce PriorEnh (coluna à esquerda) e UNetEnh (coluna à direita) para o DeepPrint (duas primeiras linhas) e FLARE com alinhamento por votação (duas últimas linhas) nos subconjuntos DB1-A e DB2-A do FVC2000.}\label{fig:histograms_enhancement}
\end{figure}


